%======================================================================
%  IEEE Conference Survey Paper  --  Template
%  Compile with:  pdflatex survey_template.tex  (run twice for refs)
%  Requires the IEEEtran class and the newtx font package (both ship
%  with a full MiKTeX / TeX Live installation).
%======================================================================
\documentclass[conference]{IEEEtran}
\IEEEoverridecommandlockouts

% ------------------------- Typography ---------------------------------
% newtx gives a clean, professional Times New Roman look for both text
% and mathematics -- noticeably crisper than the legacy `times' package.
\usepackage[T1]{fontenc}
\usepackage[utf8]{inputenc}
\usepackage{newtxtext}      % Times-like text font
\usepackage{newtxmath}      % matching math font
\usepackage{microtype}      % subtle kerning/spacing polish

% ------------------------- Packages -----------------------------------
\usepackage{cite}
\usepackage{amsmath,amssymb,amsfonts}
\usepackage{algorithmic}
\usepackage{graphicx}
\usepackage{textcomp}
\usepackage{xcolor}
\usepackage{booktabs}       % professional horizontal rules for tables
\usepackage{array}
\usepackage{url}
\usepackage[hidelinks]{hyperref}

\begin{document}

%======================================================================
% TITLE
% (24pt, Centered, Title Case -- clearly identifies the paper as a survey)
%======================================================================
\title{A Comprehensive Survey on AI-Assisted Speech Therapy for Children with Autism Spectrum Disorder\\
{\footnotesize \textsuperscript{}}
\thanks{Note: Sub-titles are not captured in Xplore and should not be used.}
}

%======================================================================
% AUTHORS
% Formatted per the official IEEE conference template (conference_101719):
% ordinal-numbered blocks arranged in a 3 x 2 grid.
% (10pt, Centered)
%======================================================================
\newcommand{\authcol}[1]{%
  \begin{minipage}[t]{0.31\textwidth}\centering\small #1\end{minipage}}

\author{%
% ---------- Row 1 ----------
\authcol{%
  \textbf{1\textsuperscript{st} Rohith Kumar Dhamagatla}\\[2pt]
  \textit{Dept. of CSE}\\
  \textit{Amrita Vishwa Vidyapeetham}\\
  Coimbatore, India\\
  cb.sc.u4cse23018@cb.students.amrita.edu}
\hfill
\authcol{%
  \textbf{2\textsuperscript{nd} Thungamitta Venkataramana}\\[2pt]
  \textit{Dept. of CSE}\\
  \textit{Amrita Vishwa Vidyapeetham}\\
  Coimbatore, India\\
  cb.sc.u4cse23055@cb.students.amrita.edu}
\hfill
\authcol{%
  \textbf{3\textsuperscript{rd} Ajay Bhargava Jashwanth Reddy Vustelamuri}\\[2pt]
  \textit{Dept. of CSE}\\
  \textit{Amrita Vishwa Vidyapeetham}\\
  Coimbatore, India\\
  cb.sc.u4cse23058@cb.students.amrita.edu}
\\[12pt]
% ---------- Row 2 ----------
\authcol{%
  \textbf{4\textsuperscript{th} Challa Kalyan Kumar Reddy}\\[2pt]
  \textit{Dept. of CSE}\\
  \textit{Amrita Vishwa Vidyapeetham}\\
  Coimbatore, India\\
  cb.sc.u4cse20360@cb.students.amrita.edu}
\hfill
\authcol{%
  \textbf{5\textsuperscript{th} Sholingram Hemanth}\\[2pt]
  \textit{Dept. of CSE}\\
  \textit{Amrita Vishwa Vidyapeetham}\\
  Coimbatore, India\\
  cb.sc.u4cse23504@cb.students.amrita.edu}
\hfill
\authcol{%
  \textbf{6\textsuperscript{th} Dr. Venkataraman D}\\[2pt]
  \textit{Associate Professor}\\
  \textit{Dept. of CSE}\\
  \textit{Amrita Vishwa Vidyapeetham}\\
  Coimbatore, India\\
  d\_venkataraman@cb.amrita.edu}
}

\maketitle

%======================================================================
% ABSTRACT  (Bold, Fully Justified, 150-250 words)
%======================================================================
\begin{abstract}
Autism Spectrum Disorder (ASD) is a neurodevelopmental condition in which speech, language, and social-communication deficits are core clinical features, yet access to consistent, individualized speech therapy remains constrained by cost, clinician scarcity, and geography. The convergence of self-supervised speech representation learning, multimodal artificial intelligence (AI), and gamified mobile delivery has, over the last five years, opened a path toward scalable, engaging, and objective computer-assisted therapy for autistic children. This paper presents a comprehensive survey of twenty-five studies published between 2013 and 2026 that collectively define this emerging pipeline. We motivate the survey by the absence of a unifying taxonomy that connects low-level acoustic modelling to deployable, child-facing therapy systems. The literature is organized into six themes: (i) self-supervised speech representations (wav2vec2.0) for child and atypical speech; (ii) acoustic and prosodic biomarkers of ASD; (iii) multimodal AI for early screening and personalized intervention; (iv) digital games, mobile apps, and interactive therapy platforms; (v) speech emotion and affect recognition; and (vi) enabling technologies for enhancement, forced alignment, and privacy-preserving edge deployment. For each theme we compare methods, datasets, reported performance, and limitations. Our synthesis reveals a persistent adult-to-child domain gap, a scarcity of large open child-speech corpora, weak clinical-efficacy evidence for consumer apps, and limited model explainability. We conclude by identifying research directions---syllable-level pronunciation assessment, child-voice normalization, federated personalization, and standardized benchmarks---needed to build trustworthy, world-class therapy platforms.
\end{abstract}

%======================================================================
% INDEX TERMS  (3-5 keywords)
%======================================================================
\begin{IEEEkeywords}
Speech Therapy, Autism Spectrum Disorder, Pronunciation Assessment, Deep Learning, Gamification, Survey
\end{IEEEkeywords}

%======================================================================
% I. INTRODUCTION
%======================================================================
\section{Introduction}
\IEEEPARstart{A}{utism} Spectrum Disorder (ASD) is a lifelong neurodevelopmental disorder defined by persistent deficits in social communication and by restricted, repetitive patterns of behaviour. Its prevalence has risen sharply worldwide, and a substantial fraction of affected children are minimally verbal or non-verbal, making expressive speech and language among the most consequential therapeutic targets. Speech production is a coordinated process spanning respiratory, phonatory, articulatory, and cognitive subsystems; in ASD these subsystems frequently exhibit atypical prosody, imprecise articulation, and irregular voice quality. Early, intensive, and individualized intervention is widely regarded as the single most effective determinant of long-term communicative outcome. In practice, however, such intervention depends on the scarce and costly time of speech-language pathologists (SLPs), producing long waiting lists and highly uneven access across regions and income levels.

\textit{Background.} The last decade has transformed the computational tools available for this problem. Self-supervised models such as wav2vec2.0 learn powerful speech representations from unlabelled audio; deep and multimodal architectures fuse acoustic, linguistic, and physiological cues; gamified mobile applications deliver practice at home; and edge and federated learning make private, on-device inference feasible. Each advance has been demonstrated in isolation, but a therapy platform for autistic children must integrate all of them---accurate child-speech recognition, objective pronunciation scoring, affect-aware engagement, and privacy-preserving deployment---into a single trustworthy system.

\textit{Motivation.} Existing surveys tend to be narrow: they address emotion recognition, digital games, or disorder classification separately, and are frequently confined to a single modality or clinical sub-problem. None provides a taxonomy that traces the full path from low-level acoustic modelling to a deployed, child-facing, gamified therapy application. Given the rapid 2023--2026 wave of self-supervised and generative-AI results, an integrative and up-to-date synthesis is timely and necessary.

\textit{Scope.} We review twenty-five peer-reviewed studies published between 2013 and 2026 that intersect speech technology and childhood communication disorders, with an emphasis on ASD. We include foundational enabling technologies (representation learning, alignment, enhancement, edge and federated learning) where they are directly relevant to a therapy pipeline, and we exclude work on adult-only pathologies, non-speech interventions, and purely hardware robotics that fall outside the software-therapy scope.

\textit{Contributions.} The primary contributions of this paper are:
\begin{itemize}
    \item A six-theme taxonomy that, for the first time, unifies self-supervised child-speech modelling, ASD biomarkers, multimodal screening, interactive therapy delivery, affect recognition, and enabling infrastructure.
    \item A critical, study-level comparison of methods, datasets, reported performance, and limitations across all twenty-five works, consolidated in Table~\ref{tab:comparison}.
    \item An identification of the cross-cutting gaps---the adult--child domain mismatch, data scarcity, weak efficacy evidence, and limited explainability---that currently block clinical-grade deployment.
    \item A concrete research agenda spanning syllable-level pronunciation assessment, child-voice normalization, federated personalization, and standardized benchmarks.
\end{itemize}

\textit{Organization.} The remainder of this paper is organized as follows: Section~\ref{sec:background} presents the necessary background and preliminaries; Section~\ref{sec:survey} develops the literature taxonomy; Section~\ref{sec:analysis} provides a comparative analysis; Section~\ref{sec:future} discusses open issues and future directions; and Section~\ref{sec:conclusion} concludes the survey.

%======================================================================
% II. BACKGROUND / PRELIMINARIES
%======================================================================
\section{Background / Preliminaries}
\label{sec:background}
To make the survey self-contained, this section defines the core concepts, signal representations, and learning paradigms that recur throughout the reviewed literature.

\subsection{Speech Production and ASD Speech Characteristics}
Speech is generated by the coordinated action of the respiratory (airflow), phonatory (vocal-fold vibration), and articulatory (tongue, lips, jaw) subsystems, under continuous cognitive control. Children with ASD frequently present atypical \emph{prosody}---the suprasegmental intonation, rhythm, loudness, and voice quality of speech---together with imprecise vowel articulation and irregular timing. These deviations are clinically salient and, importantly for this survey, acoustically measurable, which is what makes automated assessment feasible.

\subsection{Acoustic Features}
Classical pipelines convert the waveform into compact features. \emph{Mel-Frequency Cepstral Coefficients} (MFCCs) summarize the short-term spectral envelope on a perceptual (mel) scale and remain a strong baseline. Prosodic descriptors include the fundamental frequency $F_0$, jitter and shimmer (cycle-to-cycle variation of frequency and amplitude), and energy contours. Articulatory descriptors such as \emph{Vowel Space Area} (VSA), computed from the first two formants $F_1$ and $F_2$, quantify articulatory precision. Feature sets such as eGeMAPS standardize these measurements for reproducibility.

\subsection{Self-Supervised Representation Learning}
Modern systems replace hand-crafted features with representations learned from raw audio. \emph{wav2vec2.0} encodes the waveform with a convolutional feature extractor, masks spans of the latent sequence, and solves a contrastive task against quantized targets, learning phonetically rich embeddings from unlabelled speech. A pretrained encoder is then \emph{fine-tuned} on a small labelled set with a Connectionist Temporal Classification (CTC) objective for recognition, or with a classification head for detection tasks. Because pretraining is data-hungry and child speech is scarce, transfer learning from adult corpora is the norm---and the source of the adult-to-child domain gap analysed later.

\subsection{Alignment, Enhancement, and Pronunciation Scoring}
\emph{Forced alignment} maps a known phoneme sequence to time intervals in the audio, typically via dynamic time warping (DTW) or the CTC posteriorgram, and underpins per-syllable feedback. \emph{Goodness of Pronunciation} (GOP) scores each phone from posterior probabilities to estimate pronunciation quality. \emph{Vocal-tract length normalization} (VTLN) warps the frequency axis to compensate for the shorter vocal tract of children, reducing mismatch with adult-trained models. \emph{Speech enhancement} and neural \emph{vocoders} improve or reconstruct signal quality for noisy, low-resource inputs.

\subsection{Deployment Paradigms}
Two paradigms recur when moving from laboratory to clinic. \emph{Edge/TinyML} deploys quantized models on microcontrollers or phones for low-latency, offline, privacy-preserving inference. \emph{Federated learning} trains a shared global model across many devices without centralizing raw audio, exchanging only model updates---attractive for sensitive child-health data, at a measurable privacy--accuracy trade-off. \emph{Gamification} and reinforcement-learning-based adaptivity provide the engagement and personalization layer essential for young autistic users.

%======================================================================
% III. LITERATURE SURVEY / TAXONOMY
%======================================================================
\section{Literature Survey / Taxonomy}
\label{sec:survey}
The reviewed literature is organized into six thematic categories rather than presented chronologically, so that methodological lineages and their trade-offs become explicit. The taxonomy deliberately traverses the full engineering pipeline of an AI-assisted speech-therapy system for children with Autism Spectrum Disorder (ASD): it begins with the \emph{representation-learning} substrate that turns raw child speech into machine-usable features; proceeds to the \emph{biomarkers} that give those features clinical meaning; ascends to \emph{multimodal screening and intervention} systems that act on them; then to the \emph{interactive delivery} layer---games, apps, and feedback interfaces---through which therapy actually reaches the child; incorporates the \emph{affective} channel that sustains a young user's engagement; and closes with the \emph{enabling infrastructure} of enhancement, alignment, and privacy-preserving deployment on which any trustworthy platform must rest. Table~\ref{tab:comparison} complements the narrative with a compact, study-level comparison. For each work we consistently report the problem addressed, the method and model family, the dataset and evaluation protocol, the headline quantitative result, and the principal limitation, so that the strengths and weaknesses of each contribution can be weighed on common terms. Within every category the discussion moves from the most foundational or highest-performing contribution toward the more specialized or application-specific ones, and each category ends with a short synthesis that positions its papers against one another and against the survey's overarching goal.

\subsection{Self-Supervised Speech Representations for Child and Atypical Speech}
The first and most foundational category comprises studies that adapt large self-supervised learning (SSL) models---principally \emph{wav2vec2.0}---to the low-resource, high-variability regime of child and pathological speech. The shared premise is that representations learned from massive unlabelled audio can be transferred, with modest labelled data, to tasks for which annotated child corpora are prohibitively scarce.

Jain \textit{et al.}~\cite{ref1} provide the most systematic treatment of this premise for child Automatic Speech Recognition (ASR). They exhaustively varied the pretraining and finetuning data---adult speech, child speech, and mixtures---for wav2vec2 and measured the point at which child data begins to dominate performance. Finetuning with as little as ten hours of child speech already surpassed the unmodified BASE~960 baseline, and their best configurations achieved word-error rates (WER) of 7.42 on MyST, 2.91 on PF-STAR, and 12.77 on CMU~Kids using cleaned dataset variants that the authors released for reproducibility. The study's strength is its careful data-ablation methodology and open resources; its limitation is that even the best model still degrades on the youngest and most spontaneous speakers, underscoring the adult-to-child domain gap that recurs throughout this survey.

Lee \textit{et al.}~\cite{ref2} shift the objective from transcription to \emph{detection}, proposing an end-to-end network that classifies infants as ASD or typically developing (TD) directly from voice. They combined two complementary feature extractors---a bottleneck autoencoder over eGeMAPS descriptors and the context vector of a pretrained wav2vec2.0 encoder applied to the raw waveform---feeding a bidirectional LSTM classifier. Evaluated on recordings collected at Seoul National University Bundang Hospital and a Living Lab, the joint-optimized wav2vec2.0 model raised accuracy from 64.74\% to 71.66\% and unweighted average recall (UAR) from 65.04\% to 70.81\% over the autoencoder baseline. The work demonstrates that SSL embeddings encode ASD-discriminative cues without hand-crafted features, though the modest absolute accuracy reflects the difficulty of infant data and small cohorts.

Al~Futaisi \textit{et al.}~\cite{ref3}, in the Noor Project, foreground \emph{fairness} and data efficiency. Facing a shortage of labelled ASD speech, they compared two finetuning strategies: Discriminative Fine-Tuning (D-FT), pretrained on a related dataset before task adaptation, and wav2vec2.0 Fine-Tuning (W2V2-FT), which leverages self-supervised representations from a larger unrelated corpus. On a binary typicality task (TD vs.\ atypical), D-FT reached 94.8\% UAR versus 91.5\% for W2V2-FT; on a harder four-class diagnosis task (TD, ASD, dysphasia, PDD-NOS) D-FT still led at 60.9\% versus 54.3\% UAR. The sharp drop from binary to multi-class quantifies how limited labelled data caps fine-grained diagnosis---a central obstacle for downstream therapy targeting.

Javanmardi \textit{et al.}~\cite{ref4} extend the SSL paradigm to a neighbouring pathology, dysarthria, whose atypical prosody and imprecise articulation overlap phenomenologically with ASD speech. Using the pre-trained wav2vec model as a frozen feature extractor over the UA-Speech database, they analysed which layers best serve which task. Embeddings from the \emph{first} layer gave the best detection performance (a 1.23\% absolute accuracy gain over the spectrogram baseline), whereas embeddings from the \emph{final} layer were optimal for severity-level classification (a 10.62\% absolute gain over MFCCs). This layer-wise insight---that shallow layers capture presence and deep layers capture degree---is directly transferable to graded pronunciation assessment in therapy.

Khurana \textit{et al.}~\cite{ref5} generalize SSL transfer across languages, which matters for multilingual and low-resource therapy settings. Their three-step cross-lingual framework inserts a semantic knowledge-distillation stage into the standard two-step XLS-R pipeline, explicitly targeting the large cross-lingual transfer gap (TRFGap) between high- and low-resource languages. On the CoVoST-2 speech-translation benchmark the method delivered significant BLEU gains for low-resource languages and a measurable TRFGap reduction. Although framed around translation, the underlying lesson---that an intermediate adaptation step can bridge a representation mismatch---is precisely what child-speech systems need to cross the adult-to-child gap.

\emph{Synthesis.} Across~\cite{ref1,ref2,ref3,ref4,ref5}, self-supervised pretraining is established as the dominant strategy for scarce-data speech, and wav2vec2.0 as its de-facto backbone. The recurring themes---data-efficient finetuning, layer-wise task specialization, and intermediate adaptation to bridge domain gaps---form the technical foundation on which the remaining categories build, but they also expose the persistent adult-to-child mismatch and the ceiling that limited labels impose on fine-grained diagnosis.

\subsection{Acoustic and Prosodic Biomarkers of ASD}
Where the first category asks \emph{how} to represent child speech, the second asks \emph{what} in that speech signals ASD. These studies treat the acoustic signal itself as a diagnostic biomarker, isolating prosodic, articulatory, and motor-coordination features that distinguish autistic from typically developing children and that, in a therapy context, define the very targets a system must measure and improve.

Talkar \textit{et al.}~\cite{ref6} advance a motor-coordination hypothesis: that speech-production errors in ASD arise partly from deficits in coordinating the articulatory, laryngeal, and respiratory subsystems, deficits also visible in fine motor control. In a pilot of five children with ASD and five neurotypical peers, they computed correlations across acoustic, video, and handwriting time series during speech and handwriting tasks, and used eigenvalues of these correlation structures as proxies for coordination complexity. The derived features discriminated the two groups and exposed differences in cross-subsystem coupling. The study is conceptually important for linking speech to broader motor control, but its very small cohort limits statistical generalization---emblematic of the data scarcity pervading this category.

Godel \textit{et al.}~\cite{ref7} scale the biomarker idea to prosody in a substantially larger sample of 74 autistic preschoolers. They first built an original diarization pipeline to extract children's vocalizations from naturalistic social-interaction recordings, then derived voice-quality and prosodic metrics. The analysis revealed a robust prosodic signature of developmental level and ASD symptom severity; strikingly, some measures predicted communicative outcome one year later in children who had not yet acquired speech. This work establishes prosody as a quantifiable, even prognostic, biomarker and demonstrates that automated diarization can operate on ecologically valid recordings---both prerequisites for unobtrusive, longitudinal therapy monitoring.

Sandoval \textit{et al.}~\cite{ref8} contribute at the segmental level with an automatic, formant-based estimator of \emph{Vowel Space Area} (VSA), the area spanned by vowels in the $F_1$--$F_2$ plane and a classical index of articulatory precision. By removing the subjectivity and labour of manual formant measurement, their method makes VSA usable at scale as an objective articulation metric. Although validated in the context of dysarthric and Parkinsonian speech, the technique transfers directly to ASD, where imprecise vowel articulation is common, and it exemplifies the kind of interpretable, clinically grounded feature that complements opaque deep embeddings.

Abinath \textit{et al.}~\cite{ref9} examine the speaker-identity dimension of ASD speech. Motivated by evidence that individuals with ASD classify voices by exact acoustic characteristics rather than qualitative speaker traits, they investigated CNN-learned feature representations for characterizing and classifying the voices of talkative children with ASD, arguing that child vocalizations can serve as an objective developmental marker. Their pipeline replaces hand-crafted features with learned ones to improve robustness, though accuracy remains bounded by the limited size of available training data---once again the recurring constraint.

\emph{Synthesis.} Together~\cite{ref6,ref7,ref8,ref9} confirm that ASD leaves measurable traces across motor coordination, prosody, vowel articulation, and speaker-specific acoustics. They supply the interpretable clinical targets that the SSL models of the previous category can learn to detect, and the prognostic findings of~\cite{ref7} in particular hint at therapy systems that track development over time. The shared weakness is cohort size and heterogeneity, which motivates the larger-scale, multimodal screening systems considered next.

\subsection{Multimodal AI for Early Screening and Personalized Intervention}
The third category scales isolated biomarkers into deployable, decision-making systems, and marks the field's transition from passive measurement toward active, closed-loop intervention. Its two exemplars are notable for their scale and for coupling prediction with action.

Bae \textit{et al.}~\cite{ref10} address early screening at population scale with a two-stage multimodal framework trained on data from 1242 children aged 18--48 months, collected through a mobile application capturing parent--child interaction audio alongside standardized screening instruments (M-CHAT, SCQ-L, SRS). Rather than using only questionnaire scores, they apply natural-language processing to the \emph{text} of the screening items to extract behavioural descriptors, which they fuse with audio features. Stage~1 separates typically developing from high-risk/ASD children with AUROC 0.942; Stage~2 distinguishes high-risk from ASD by combining task-success data with SRS text, reaching AUROC 0.914 and 85.2\% accuracy. Predicted risk categories agreed with gold-standard ADOS-2 assessments in 79.59\% of cases and correlated strongly ($r=0.830$, $p<0.001$). The work is a compelling demonstration that mobile-collected, multimodal data can approach expert-level screening, though its single-programme cohort leaves external validity open.

Zheng \textit{et al.}~\cite{ref11} complete the loop from diagnosis to \emph{personalized intervention}. Their system pairs a deep neural network---trained across datasets spanning toddlers to adolescents and achieving 96.98\% accuracy, 97.65\% precision, 96.74\% recall, and 99.75\% ROC-AUC, outperforming Random Forest and Logistic Regression baselines---with a Deep Deterministic Policy Gradient (DDPG) reinforcement-learning agent. Using multi-strategy feature selection (LASSO, Random Forest) to identify salient predictors such as the Q-CHAT-10 score and ethnicity, the DDPG agent simulated personalized intervention strategies over twelve monthly cycles, optimizing intervention type, frequency, and intensity. The simulation reported up to 25\% improvement in social skills, 30\% reduction in behavioural issues, and 20\% gain in emotional stability, with high-risk cases falling from 65\% to 25\%. The framework is a forward-looking blueprint for adaptive therapy, but its intervention results rest on \emph{virtual} data, so real-world efficacy remains to be established.

\emph{Synthesis.} Papers~\cite{ref10} and~\cite{ref11} embody the category's defining shift: from single-shot classification toward scalable screening and closed-loop, personalized intervention. Both leverage mobile data collection and multimodal fusion, and both point toward the reinforcement-learning-driven adaptivity that a world-class therapy platform requires. Their common caveat---reliance on single-site or simulated data---frames the trust and validation challenges revisited in Section~\ref{sec:future}.

\subsection{Digital Games, Mobile Apps, and Interactive Therapy Platforms}
The fourth category is the application layer through which therapy actually reaches the child. It contains both systematic reviews that map the landscape of consumer tools and system papers that build concrete interactive platforms, together defining the gamification, feedback, and engagement mechanisms central to any child-facing product.

Two systematic reviews frame the state of practice. Saeedi \textit{et al.}~\cite{ref12} searched four databases and, after full-text screening, analysed 27 digital games developed to treat childhood speech disorders. They found that MFCC features and the PocketSphinx engine were the most common recognition stack, that 48\% of games were designed for parent-assisted practice, and that games measurably improved children's satisfaction, motivation, and attention. Crucially, they also catalogued the chief barriers---frustration and low self-esteem after repeated failures, environmental noise, mismatched difficulty levels, and, above all, brittle speech recognition---which read as a direct requirements specification for a better system. Furlong \textit{et al.}~\cite{ref13} deliver a sobering complementary result: appraising over 5000 mobile speech-therapy apps with the Mobile Application Rating Scale, they found that fewer than 3\% even met criteria for evaluation and none were of excellent quality. Their study exposes a stark gap between the sheer availability of apps and their therapeutic validity, and proposes a reproducible method for identifying suitable tools.

On the systems side, three papers illustrate increasingly sophisticated delivery. Ortiz~Castellanos \textit{et al.}~\cite{ref14} designed a home-rehabilitation software system for ASD that integrates TensorFlow-based image classification, ARKit text-to-speech, natural-language processing, and cloud storage across mobile operating systems, and detailed how these components interoperate to let patients perform therapy anywhere while transmitting data to a clinician. Chueh \textit{et al.}~\cite{ref15} move from architecture to evidence, evaluating a generative-AI-powered tablet platform for home-based clinical language therapy in a single-arm pre/post pilot of 22 children aged 2--12 (19 completing both phases); 2000 audio recordings were collected and participants stratified by cumulative AI-usage hours, providing early feasibility evidence that generative AI can extend therapy from clinic to home. Katz \textit{et al.}~\cite{ref16} represent the high-fidelity feedback frontier with Opti-Speech, a real-time 3D visual-feedback system that tracks tongue and jaw movement via electromagnetic articulography and renders an avatar of the articulators inside a transparent head, giving talkers augmented feedback when they reach correct places of articulation; while powerful, its reliance on specialized magnetometer hardware makes it clinic-bound rather than home-deployable.

Complementing software approaches, Aftab \textit{et al.}~\cite{ref17} review augmentative and alternative communication (AAC) for minimally verbal children with ASD. Screening 468 articles down to eight rigorous studies and applying Tau-U and improvement-rate-difference analyses (0.80--1.00 for single-case designs, 0.90--0.95 for RCTs), they conclude that AAC aids effectively increase communication and that high-tech aids---which children also prefer---outperform low-tech ones for social communication, interaction, and speech production. This grounds the case for technology-rich, engaging interfaces in controlled evidence.

\emph{Synthesis.} Category four supplies the engagement and delivery layer~\cite{ref12,ref13,ref14,ref15,ref16,ref17} that turns algorithms into therapy children will actually use, and the AAC evidence~\cite{ref17} confirms that high-tech, preferred interfaces yield real communicative gains. Yet the two large reviews~\cite{ref12,ref13} converge on a critical warning that echoes across the survey: engagement mechanisms are abundant, but robust child-speech recognition and rigorous efficacy evidence remain scarce---precisely the deficits that the SSL, biomarker, and enabling-technology categories must jointly remedy.

\subsection{Speech Emotion and Affect Recognition}
The fifth category addresses the affective channel, which is indispensable for a technology intended to engage autistic children: a therapy system must sense frustration, boredom, or delight to adapt its behaviour, and emotional expression is itself frequently atypical in ASD. Two complementary reviews define the methodological landscape.

Nguyen \textit{et al.}~\cite{ref18} provide a comprehensive review of \emph{multimodal fusion} for speech emotion recognition (SER) synthesizing publications from 2020 to 2024. They argue that unimodal SER, though simple and efficient, falters under noise and ambiguity, motivating multimodal integration, and they taxonomize the principal fusion strategies---early, late, deep, and hybrid fusion---alongside data-representation and data-translation techniques, attention mechanisms, and graph-based fusion. Assessing these across standard SER datasets, they map performance against challenges of data alignment, noise management, and computational demand. This fusion taxonomy is directly reusable as the design vocabulary for combining a child's voice, video, and interaction signals in a therapy loop.

Landowska \textit{et al.}~\cite{ref19} narrow the lens precisely to the target population with a PRISMA-guided systematic review of automatic emotion recognition in children with autism. Surveying diverse observation channels---facial expression, speech prosody, and physiological signals---they find that basic emotions (especially happiness, fear, and sadness) dominate the recognized set, that single-channel approaches are more common than multichannel, and that where fusion is used, early fusion prevails, with SVMs and neural networks the most popular classifiers. They further surface qualitative issues around participant-group construction and, critically, the scarcity of open, well-labelled datasets and the transient unavailability of individual channels.

\emph{Synthesis.} Papers~\cite{ref18} and~\cite{ref19} establish that affect recognition can supply the real-time feedback signal an adaptive therapy system needs, and the fusion strategies catalogued in~\cite{ref18} align neatly with the multimodal screening methods of Category three. Once more, however, the binding constraint is data: the shortage of open, richly labelled emotional corpora for autistic children, flagged explicitly in~\cite{ref19}, mirrors the scarcity seen for child ASR and biomarker studies.

\subsection{Enabling Technologies: Enhancement, Alignment, and Edge/Privacy Deployment}
The sixth and final category collects the infrastructure that renders every preceding system practical, trustworthy, and deployable. Its papers span a broad survey of the field, signal enhancement, fine-grained alignment for pronunciation scoring, on-device inference, and privacy-preserving training---the substrate on which a real therapy platform is engineered.

Remya \textit{et al.}~\cite{ref20} offer a wide-angle systematic review that situates the entire domain: synthesizing 231 studies (2007--2024) on machine- and deep-learning methods for classifying and enhancing sixteen distinct speech and language disorders. They chart the progression from traditional classifiers such as support vector machines to convolutional and recurrent networks and hybrid multimodal models, analyse acoustic, prosodic, and linguistic feature families, and review deep enhancement techniques---multi-scale feature fusion, noise-aware variational autoencoders, LLM-integrated communication aids, and ultrasound-guided articulation reconstruction---for improving intelligibility. Critically, they identify the field's binding limitations: a lack of longitudinal and cross-cultural datasets and the limited explainability of AI models in clinical contexts, nominating LLMs and adaptive techniques as routes to disorder-agnostic, personalized care.

For signal quality, Govalkar \textit{et al.}~\cite{ref21} present a fair comparison of recent neural vocoders in a speech-reconstruction setting, resynthesizing waveforms from mel-scaled spectrograms and evaluating six neural vocoders against two classical phase-reconstruction methods using the MUSHRA listening-test standard. By weighing perceptual quality against computational requirements, their study offers practitioners a concrete guideline for choosing enhancement and synthesis components---relevant wherever a therapy system must clean or regenerate noisy, low-resource child audio.

For fine-grained feedback, Zhu \textit{et al.}~\cite{ref22} tackle alignment, the technical prerequisite for per-syllable pronunciation scoring. They curated IPAPACK, a massively multilingual phonemically transcribed corpus spanning 115 languages, and proposed CLAP-IPA, a phoneme-speech contrastive embedding model enabling open-vocabulary matching between arbitrary speech and phonemic sequences; a neural forced aligner, IPA-ALIGNER, further finetunes it for phone-to-audio alignment. Tested on 95 unseen languages, the model generalizes without adaptation, and temporal alignments emerge from contrastive training---capabilities that make language-agnostic keyword spotting and forced alignment, and therefore multilingual pronunciation feedback, feasible.

For deployment, two papers address the practical realities of running such systems at the edge and under privacy constraints. Barovic and Moin~\cite{ref23} train and deploy a quantized 1D convolutional network for speech recognition on a highly resource-constrained microcontroller (Arduino Nano~33 BLE Sense), achieving up to 97\% accuracy on a purpose-built dataset while handling 23 distinct keywords---evidence that responsive, offline, privacy-preserving inference is achievable on TinyML hardware suited to ambient assisted living and assistive applications. Woubie and B\"ackstr\"om~\cite{ref24} address the training side, investigating federated learning for speaker recognition with and without secure aggregation, using a generative adversarial network to synthesize imposter data at the edge so that raw audio never leaves the device. On VoxCeleb-1 they show that federated training preserves privacy and can even improve average equal error rate over individual models, while quantifying the inherent privacy--accuracy trade-off---directly relevant to sensitive child-health data.

\emph{Synthesis.} This category~\cite{ref20,ref21,ref22,ref23,ref24} provides the connective tissue of the whole pipeline: a field-level map and explainability agenda~\cite{ref20}, enhancement and synthesis components~\cite{ref21}, the alignment machinery for syllable-level feedback~\cite{ref22}, and the edge~\cite{ref23} and federated~\cite{ref24} mechanisms that make private, offline, personalized deployment possible. Read together with the preceding five categories, these enabling technologies complete the blueprint for a trustworthy, scalable, and engaging AI-assisted speech-therapy platform---while their shared emphasis on explainability, data diversity, and privacy trade-offs foreshadows the open challenges analysed in the remainder of the paper.

%======================================================================
% IV. COMPARATIVE ANALYSIS & DISCUSSION
%======================================================================
\section{Comparative Analysis and Discussion}
\label{sec:analysis}
Instead of just summarizing papers, a strong IEEE survey critically analyzes them. This section frequently utilizes tables to compare multiple studies at a glance.

\begin{table}[!t]
\renewcommand{\arraystretch}{1.3}
\caption{Comparative Analysis of Existing Literature}
\label{tab:comparison}
\centering
\footnotesize
\begin{tabular}{@{}c c p{1.3cm} p{1.5cm} p{1.5cm}@{}}
\toprule
\textbf{Ref.} & \textbf{Year} & \textbf{Method / Focus} & \textbf{Key Contribution} & \textbf{Limitation}\\
\midrule
\cite{ref1}  & 2023 & wav2vec2 child ASR & Low WER with $\leq$10\,h child data & Adult--child domain gap\\
\cite{ref3}  & 2025 & Transfer learning (ASD) & 94.8\% UAR typicality task & Small labelled corpora\\
\cite{ref10} & 2025 & Multimodal screening & AUROC 0.942, ADOS-2 agreement & Single-site data\\
\cite{ref11} & 2025 & DNN + DDPG & Adaptive personalized therapy & Virtual/simulated data\\
\cite{ref16} & 2014 & EMA visual feedback & Real-time articulatory avatar & Intrusive hardware\\
\cite{ref24} & 2021 & Federated learning & Privacy-preserving training & Privacy--accuracy trade-off\\
\bottomrule
\end{tabular}
\end{table}

Table~\ref{tab:comparison} and the taxonomy together reveal several cross-cutting patterns. First, there is a clear \emph{methodological migration} from hand-crafted acoustic features and classical classifiers (SVMs, MFCC pipelines) toward self-supervised transformers: the strongest recent results for child and atypical speech~\cite{ref1,ref2,ref3,ref4} all build on wav2vec2.0 rather than bespoke feature engineering. Second, a \emph{modality-fusion trend} is evident---the highest screening performance is obtained by combining voice with questionnaire text or task-success signals~\cite{ref10}, echoing the fusion taxonomy formalized for emotion recognition~\cite{ref18,ref19}. Third, the field is beginning to close the loop: rather than stopping at classification, recent systems couple prediction with \emph{adaptive intervention}~\cite{ref11} and generative, interactive delivery~\cite{ref14,ref15}.

Alongside these advances, consistent \emph{trade-offs} recur. Performance is repeatedly bought with data that does not transfer: models tuned on adult corpora degrade on children~\cite{ref1}, and impressive laboratory accuracies rest on small, single-site, or simulated cohorts~\cite{ref6,ref10,ref11}. Deployment introduces a second trade-off axis---edge quantization and federated aggregation buy privacy and offline operation at the cost of accuracy~\cite{ref23,ref24}. Finally, the two large app/game reviews~\cite{ref12,ref13} expose a maturity gap: engagement mechanisms are plentiful, but rigorous efficacy evidence and robust child-speech recognition remain rare. Chronologically, the 2013--2020 works emphasize biomarkers and articulatory feedback~\cite{ref8,ref16,ref6}, whereas the 2022--2026 works pivot decisively toward self-supervised, multimodal, and generative-AI systems~\cite{ref3,ref10,ref11,ref15}, confirming that the domain is in a rapid transition phase that justifies this survey.

%======================================================================
% V. OPEN ISSUES & FUTURE RESEARCH DIRECTIONS
%======================================================================
\section{Open Issues and Future Research Directions}
\label{sec:future}
Synthesizing the limitations identified above, we distil five open challenges and the directions most likely to resolve them.

\subsection{The Adult-to-Child Domain Gap}
State-of-the-art recognizers are pretrained on abundant adult speech and degrade on the shorter vocal tracts, higher pitch, and greater variability of children~\cite{ref1,ref2}. Vocal-tract length normalization, child-specific data augmentation, and self-supervised pretraining on unlabelled child audio are promising, but a principled, reproducible child-normalization pipeline integrated with modern encoders is still missing.

\subsection{Data Scarcity and the Absence of Open Benchmarks}
Nearly every reviewed study laments small, private, or single-site datasets~\cite{ref3,ref6,ref10,ref19}. The lack of large, ethically shared, demographically diverse child-ASD speech corpora---and of standardized evaluation protocols---prevents fair comparison and reproducibility. Community benchmarks, privacy-preserving data-sharing via federated learning~\cite{ref24}, and synthetic augmentation are needed.

\subsection{From Detection to Fine-Grained, Actionable Feedback}
Most systems output a coarse label (ASD/TD, severity level) rather than the syllable- or phoneme-level guidance a therapy loop requires. Forced alignment and Goodness-of-Pronunciation scoring~\cite{ref22} point toward per-syllable feedback, but their integration into engaging, child-facing therapy remains largely unexplored.

\subsection{Weak Clinical Efficacy and Trust}
Consumer apps are abundant yet rarely validated: fewer than 3\% met quality criteria in one large appraisal~\cite{ref13}, and games often lack controlled efficacy trials~\cite{ref12}. Coupled with the limited explainability of deep models in clinical use~\cite{ref20}, this undermines clinician and parent trust. Explainable, clinically co-designed, and rigorously trialled systems are a priority.

\subsection{Privacy-Preserving, Affect-Aware, On-Device Personalization}
Child-health audio is highly sensitive, motivating edge and federated deployment~\cite{ref23,ref24}; effective therapy additionally demands affect-aware engagement~\cite{ref18,ref19} and reinforcement-learning-based personalization~\cite{ref11}. Unifying these---personalized, emotion-aware adaptation that never centralizes raw audio---is an open systems challenge for the next generation of platforms.

%======================================================================
% VI. CONCLUSION
%======================================================================
\section{Conclusion}
\label{sec:conclusion}
This survey examined twenty-four studies published between 2013 and 2026 at the intersection of speech technology and childhood communication disorders, with a focus on Autism Spectrum Disorder. We organized the literature into a six-theme taxonomy spanning self-supervised child-speech representations, acoustic and prosodic biomarkers, multimodal screening and personalized intervention, digital games and interactive therapy platforms, speech emotion recognition, and the enabling technologies of enhancement, alignment, and privacy-preserving edge deployment. Across these themes, the state of the art has shifted decisively from hand-crafted features and classical classifiers toward self-supervised, multimodal, and generative-AI systems, with the most recent work beginning to close the loop between assessment and adaptive intervention. Yet the same body of work exposes recurring obstacles: a persistent adult-to-child domain gap, a scarcity of large open child-speech corpora and standardized benchmarks, coarse detection outputs where fine-grained feedback is needed, weak clinical-efficacy evidence for consumer applications, and limited model explainability. Addressing these---through child-voice normalization, community benchmarks, syllable-level pronunciation assessment, explainable and clinically co-designed systems, and privacy-preserving, affect-aware personalization---is, in our assessment, the path toward trustworthy, engaging, and genuinely scalable AI-assisted speech therapy. As self-supervised and generative models continue to mature, we expect the field to move from isolated proofs of concept toward integrated platforms that meaningfully extend the reach of speech-language pathologists to the children who need them most.

%======================================================================
% REFERENCES  (8pt, numbered in order of appearance)
%======================================================================
\begin{thebibliography}{00}

\bibitem{ref1} R. Jain, A. Barcovschi, M. Y. Yiwere, D. Bigioi, P. Corcoran, and H. Cucu, ``A wav2vec2-based experimental study on self-supervised learning methods to improve child speech recognition,'' \textit{IEEE Access}, vol.~11, pp.~46938--46948, 2023.

\bibitem{ref2} J. H. Lee, G. W. Lee, G. Bong, H. J. Yoo, and H. K. Kim, ``End-to-end model-based detection of infants with autism spectrum disorder using a pretrained model,'' \textit{Sensors}, vol.~23, no.~1, Art. no.~202, 2023.

\bibitem{ref3} N. D. Al Futaisi, B. W. Schuller, F. Ringeval, and M. Pantic, ``The Noor Project: Fair transformer transfer learning for autism spectrum disorder recognition from speech,'' \textit{Front. Digit. Health}, vol.~7, Art. no.~1274675, 2025.

\bibitem{ref4} F. Javanmardi, S. Tirronen, M. Kodali, S. R. Kadiri, and P. Alku, ``Wav2vec-based detection and severity level classification of dysarthria from speech,'' in \textit{Proc. IEEE Int. Conf. Acoust., Speech Signal Process. (ICASSP)}, Rhodes Island, Greece, 2023, pp.~1--5.

\bibitem{ref5} S. Khurana, N. Dawalatabad, A. Laurent, L. Vicente, P. Gimeno, V. Mingote, and J. Glass, ``Cross-lingual transfer learning for low-resource speech translation,'' arXiv:2306.00789, 2024.

\bibitem{ref6} T. Talkar \textit{et al.}, ``Assessment of speech and fine motor coordination in children with autism spectrum disorder,'' \textit{IEEE Access}, vol.~8, pp.~127535--127545, 2020.

\bibitem{ref7} M. Godel, F. Robain, F. Journal, N. Kojovic, K. Latr\`eche, G. Dehaene-Lambertz, and M. Schaer, ``Prosodic signatures of ASD severity and developmental delay in preschoolers,'' \textit{npj Digit. Med.}, vol.~6, Art. no.~99, 2023.

\bibitem{ref8} S. Sandoval, V. Berisha, R. L. Utianski, J. M. Liss, and A. Spanias, ``Automatic assessment of vowel space area,'' \textit{J. Acoust. Soc. Am.}, vol.~134, no.~5, pp.~EL477--EL483, 2013.

\bibitem{ref9} M. Abinath, S. M. Afaal, P. Anojan, R. Vithurshan, K. Pulasinga, and W. Tissera, ``Automatic speech recognition system to analyze autism spectrum disorder in young children,'' \textit{Int. J. Eng. Manag. Res.}, vol.~12, no.~5, pp.~1--8, 2022.

\bibitem{ref10} S. Bae \textit{et al.}, ``Multimodal AI for risk stratification in autism spectrum disorder: Integrating voice and screening tools,'' \textit{npj Digit. Med.}, vol.~8, Art. no.~01914-6, 2025.

\bibitem{ref11} Y. Zheng, S. Ma, X. Liu, L. Wang, M. Chen, J. Yuan, and H. Wang, ``Deep neural networks and deep deterministic policy gradient for early ASD diagnosis and personalized intervention in children,'' \textit{Sci. Rep.}, vol.~15, Art. no.~26854, 2025.

\bibitem{ref12} S. Saeedi, H. Bouraghi, M.-S. Seifpanahi, and M. Ghazisaeedi, ``Application of digital games for speech therapy in children: A systematic review of features and challenges,'' \textit{J. Healthcare Eng.}, vol.~2022, Art. no.~4814945, 2022.

\bibitem{ref13} L. Furlong, M. Morris, T. Serry, and S. Erickson, ``Mobile apps for treatment of speech disorders in children: An evidence-based analysis of quality and efficacy,'' \textit{PLoS ONE}, vol.~13, no.~8, Art. no.~e0201513, 2018.

\bibitem{ref14} A. E. Ortiz Castellanos, C.-M. Liu, and C. Shi, ``Deep mobile linguistic therapy for patients with ASD,'' \textit{Int. J. Environ. Res. Public Health}, vol.~19, no.~19, Art. no.~12857, 2022.

\bibitem{ref15} C.-H. Chueh \textit{et al.}, ``Implementation of a generative AI-powered digital interactive platform for clinical language therapy in children with language delay: A pilot study,'' \textit{Life}, vol.~15, no.~10, Art. no.~1628, 2025.

\bibitem{ref16} W. Katz \textit{et al.}, ``Opti-Speech: A real-time, 3D visual feedback system for speech training,'' in \textit{Proc. INTERSPEECH}, Singapore, 2014, pp.~1174--1178.

\bibitem{ref17} A. Aftab, C. A. Sehgal, M. M. Noohu, and G. Jaleel, ``Clinical effectiveness of AAC intervention in minimally verbal children with ASD: A systematic review,'' \textit{NeuroRegulation}, vol.~10, no.~4, pp.~239--252, 2023.

\bibitem{ref18} N. M. Nguyen, T. T. Nguyen, P.-N. Tran, C. P. Lim, N. T. Pham, and D. N. M. Dang, ``Multimodal fusion in speech emotion recognition: A comprehensive review of methods and technologies,'' \textit{Eng. Appl. Artif. Intell.}, vol.~163, Art. no.~112624, 2026.

\bibitem{ref19} A. Landowska, A. Karpus, T. Zawadzka, B. Robins, D. E. Barkana, H. Kose, T. Zorcec, and N. Cummins, ``Automatic emotion recognition in children with autism: A systematic literature review,'' \textit{Sensors}, vol.~22, no.~4, Art. no.~1649, 2022.

\bibitem{ref20} M. S. Remya, R. Raman, R. Sankaran, V. Namboodiri, and P. Nedungadi, ``Artificial intelligence for classification and enhancement of speech and language disorders: Techniques, applications, and future directions,'' \textit{IEEE Access}, vol.~13, pp.~1--28, 2025.

\bibitem{ref21} P. Govalkar, J. Fischer, F. Zalkow, and C. Dittmar, ``A comparison of recent neural vocoders for speech signal reconstruction,'' in \textit{Proc. 10th ISCA Speech Synthesis Workshop (SSW)}, Vienna, Austria, 2019, pp.~7--12.

\bibitem{ref22} J. Zhu, C. Yang, F. Samir, and J. Islam, ``The taste of IPA: Towards open-vocabulary keyword spotting and forced alignment in any language,'' in \textit{Proc. Annu. Meeting Assoc. Comput. Linguistics (ACL)}, 2024, pp.~1--16.

\bibitem{ref23} A. Barovic and A. Moin, ``TinyML for speech recognition,'' in \textit{Proc. IEEE 49th Annu. Comput., Softw., Appl. Conf. (COMPSAC)}, 2025, pp.~1--6.

\bibitem{ref24} A. Woubie and T. B\"ackstr\"om, ``Federated learning for privacy-preserving speaker recognition,'' \textit{IEEE Access}, vol.~9, pp.~149477--149487, 2021.

\end{thebibliography}

%======================================================================
% KEY IEEE FORMATTING RULES (kept as a reminder)
% - Layout: two-column (Title and Abstract span both columns).
% - Font: Times New Roman standard (newtx used here).
% - Section headings: centered, Roman numerals, small caps.
% - Subheadings: italicized, alphabetical order (A. B. C.).
% - Citations: square brackets inside punctuation, e.g. "...studies [4], [5]."
%   Do not write "Ref. [3]" except at the start of a sentence.
% - Figures/Tables: numbered, centered, captioned. Table captions ABOVE,
%   figure captions BELOW.
%======================================================================

\end{document}
